\section{A15.1b2d: full Sub-wedge~A injectivity ($e/2 \le R < d/2$) [proof retracted]}
\label{sec:p12-b2d}

\begin{center}\framebox{\begin{minipage}{0.9\textwidth}
\textbf{Retraction of proof (2026-08-22, adversarial audit).}
The proof below asserts two $19\times 19$ typed blocks with
determinant $\pm p^{8} r (p^{2} - q^{2})^{3} F$ and a Case-B third-cell
reduction identical to the b2c triangular factor.  The 19 source
positions are \emph{named as} ``the active source pairs'' but not
enumerated in this manuscript nor in the referenced audit; and the
Case-B third-cell reduction relies on the retracted b2c proof, which
used sources outside the horizon $(0, T_{0})$.  The theorem statement
is not refuted; only the proof is retracted.
Status of \textbf{Theorem~\ref{thm:p12-b2d}}: $?[O]$ pending an
explicit horizon-legal certificate with a full enumeration of the 19
typed sources and a symbolic cancellation table, for both cells and
both $\sigma$-cases.
\end{minipage}}\end{center}

The remaining slice after (the alleged) Theorem~\ref{thm:p12-b2c} is
$e/2 \le R < d/2$ with $\sigma > R$.  This section \emph{claims} to
close it; see the retraction box above.

\subsection*{Two natural thresholds}

For the defect strip $x \in (R, \sigma)$, set
\[
\kappa := e - \delta = 2e - d, \qquad \lambda := d - \sigma.
\]
Numerically $\kappa \approx 0.0849495$.  Since
$\sigma < \varepsilon_{\max} = \tfrac12\log(5/4) < 2\delta$
(numerically $2\delta \approx 0.1178$, $\varepsilon_{\max} \approx 0.1116$),
one always has $\kappa < \lambda$:
$\lambda - \kappa = d - \sigma - (e - \delta) = 2\delta - \sigma > 0$.

The defect strip decomposes as follows.

\subsection*{Case A: $\sigma \le d/2$ — two-cell decomposition}

Here $\lambda = d - \sigma \ge \sigma$, so the strip $(R, \sigma)$
lies entirely below $\lambda$, and the only splitting point inside
$(R, \sigma)$ is $\kappa$.  Two typed cells:
\[
(R, \kappa)\quad \text{and} \quad (\kappa, \sigma).
\]

\begin{lemma}[Two 19$\times$19 elimination blocks]
\label{lem:p12-b2d-19blocks}
For every $x_{0}$ in $(R, \kappa)$ or in $(\kappa, \sigma) \cap (R, d/2)$
the active source pairs at $x_{0}$ produce a $19\times 19$ typed
linear system on visibility-stable target values.  The coefficient
matrix $M_{\pm}$ (of either cell) has determinant
\[
\det M_{\pm} \;=\; \pm\, p^{8}\, r\, (p^{2} - q^{2})^{3}\, F,
\qquad
F := 2p^{4} - 3 p^{2} q^{2} - p^{2} r^{2} + q^{4} - q^{2} r^{2}.
\]
\end{lemma}

The explicit typed enumeration and the symbolic determinant reduction
are recorded in
\texttt{P12\_REFEREE\_E2E\_A15\_1B2D\_FULL\_SUBWEDGE\_A\_2026-08-22.md}.

\begin{lemma}[$F < 0$ elementary]\label{lem:p12-b2d-F-negative}
Set $B := (q/p)^{2} = 2^{-3/2}$ and $\theta := (r/p)^{2}
= (\log 3 / \log 2)(2/3)^{3/2}$.  Then
$F / p^{4} = (2 - 3B + B^{2}) - (1 + B)\, \theta$, and
\[
\frac{2 - 3B + B^{2}}{1 + B} \;<\; \frac{4}{5} \;<\; \theta,
\]
whence $F < 0$.  Both inequalities are elementary rational identities
in $\{\log 2, \log 3\}$.
\end{lemma}

\begin{proof}
Substitute $B = 1/(2\sqrt{2})$ into the left ratio and clear the
denominator: $(2 - 3B + B^{2})/(1 + B) < 4/5 \iff 53 < 38\sqrt{2}$,
and $53^{2} = 2809 < 2888 = 38^{2}\cdot 2$.

For the right inequality: $\theta = (\log 3 / \log 2)(2/3)^{3/2}
> (3/2)(2/3)^{3/2} = \sqrt{2/3}$, using $\log 3 / \log 2 > 3/2
\iff 3 > 2\sqrt{2} \iff 9 > 8$.  Then
$\sqrt{2/3} > 4/5 \iff 2/3 > 16/25 \iff 50 > 48$.
\end{proof}

\begin{corollary}[Both 19-blocks are invertible]
\label{cor:p12-b2d-invertible}
Under Lemmas~\ref{lem:p12-b2d-19blocks} and \ref{lem:p12-b2d-F-negative}
one has $\det M_{\pm} \ne 0$.
\end{corollary}

\begin{proof}
$p, r > 0$; $p^{2} - q^{2} > 0$ by Lemma~\ref{lem:p12-b2c-detfactors};
$F < 0$ by Lemma~\ref{lem:p12-b2d-F-negative}.
\end{proof}

\begin{proposition}[Case-A defect strip cleared]\label{prop:p12-b2d-caseA}
Let $2a < T_0 < c$, $e/2 \le R < d/2$, $T < S < T_0$ with
$\sigma \le d/2$.  Then $h = 0$ a.e. on $(R, \sigma)$ and $H = 0$ a.e.
on $(0, \sigma) \cap$ (visibility-stable strip).
\end{proposition}

\begin{proof}
For $x_{0} \in (R, \kappa)$ or $(\kappa, \sigma)$,
Corollary~\ref{cor:p12-b2d-invertible} forces
$v = 0$ for the $19$-vector of visibility-values at $x_{0}$; in
particular $h(x_{0}) = h(d - x_{0}) = h(a - x_{0}) = 0$.
Substituting into the defect equation (E$_{\sigma}$) at $x_{0}$
yields $H(x_{0}) = 0$.
\end{proof}

\subsection*{Case B: $\sigma > d/2$ — three-cell decomposition}

Here $\lambda = d - \sigma < \sigma$, and $\kappa < \lambda < \sigma$,
so the strip $(R, \sigma)$ splits into three typed cells:
\[
(R, \kappa),\quad (\kappa, \lambda),\quad (\lambda, \sigma),
\]
each non-empty as parameters permit.

The first two cells are covered by
Lemma~\ref{lem:p12-b2d-19blocks} exactly as in Case~A.  The third
cell $(\lambda, \sigma)$ admits a cleaner mechanism.

\begin{lemma}[Third-cell $13$-value elimination]
\label{lem:p12-b2d-thirdcell}
For $x_{0} \in (\lambda, \sigma)$, the $13$ source points
\[
d - x_{0},\; e + x_{0},\; 2d - x_{0},\; b - x_{0},\; a + x_{0},\;
3d - x_{0},\; 3e + x_{0},\; T - x_{0},\; 3e + \delta + x_{0},
\]
\[
4e + 3\delta - x_{0},\; 4e + \delta + x_{0},\; a + b - x_{0},\; T + x_{0}
\]
all lie in the source horizon $(0, T_{0})$.  Triangular reduction of
the resulting $13$ typed relations gives
\[
\frac{2 q r}{p^{3}}\,(2p^{2} - q^{2})\, h(x_{0}) = 0,
\]
identical to the b2c triangular factor
(Eq.\,\ref{eq:o3ae-triangular-b2c}).
\end{lemma}

\begin{proposition}[Case-B defect strip cleared]\label{prop:p12-b2d-caseB}
Under the Case-B hypotheses, $h = 0$ and $H = 0$ a.e. on $(R, \sigma)$.
\end{proposition}

\begin{proof}
Cells $(R, \kappa)$ and $(\kappa, \lambda)$ are cleared by
Lemma~\ref{lem:p12-b2d-19blocks}$+$Corollary~\ref{cor:p12-b2d-invertible}
exactly as in Proposition~\ref{prop:p12-b2d-caseA}.  For the third
cell $x_{0} \in (\lambda, \sigma)$, Lemma~\ref{lem:p12-b2d-thirdcell}
and Lemma~\ref{lem:p12-b2c-detfactors} force $h(x_{0}) = 0$; the
source equation at $u = T - x_{0}$ then reads
$p\, h(a - x_{0}) - q\, h(x_{0}) = 0$, giving $h(a - x_{0}) = 0$;
finally, $h(d - x_{0})$: if $d - x_{0} \le R$ it is zero by support,
otherwise $d - x_{0} \in (d - \sigma, \sigma) = (\lambda, \sigma)$ and
the same third-cell elimination gives $h(d - x_{0}) = 0$.
Substituting into (E$_{\sigma}$) yields $H(x_{0}) = 0$.
\end{proof}

\subsection*{Main theorem for Sub-wedge~A full-tail}

\begin{theorem}[Full Sub-wedge~A injectivity; proof retracted]\label{thm:p12-b2d}
Let $2a < T_0 < c$, $e/2 \le R < d/2$, $T < S < T_0$.  Then
\(\ker L_{R,S,T_0}^{\{a,b,2a\}} = \{0\}\).
\end{theorem}

\begin{proof}[Retracted proof, kept for archival adversarial review]
If $\sigma \le R$, Theorem~\ref{thm:p12-b2b} applies (Sub-wedge~A
restricted-tail).  Otherwise $\sigma > R$; by whether $\sigma \le d/2$
or $\sigma > d/2$, Proposition~\ref{prop:p12-b2d-caseA} or
\ref{prop:p12-b2d-caseB} clears the defect strip: $h = 0$ on
$(R, \sigma)$ and $H = 0$ on $(0, \sigma)$ (visibility-stable part).

With $H = 0$ on the visibility-stable part of $(0, \sigma)$ and
$h = 0$ on $(R, \sigma)$, equation (E$_{\sigma}$) is homogeneous on
$(\sigma, a)$ and the weighted-reflection kill of A15.1b2b (Sub-wedge~A
restricted-tail, Theorem~\ref{thm:p12-b2b}) applies on that domain,
giving $h = 0$ on $(\sigma, a)$.  For $0 < t < R$ the defect equation
now involves only values in the killed region, yielding $H(t) = 0$;
the tail $(T, S)$ is thus entirely zero.  Support reduces to $(R, T)$,
and Theorem~\ref{thm:p12-b1} (A15.1b1) closes the residual.
\end{proof}

\begin{corollary}[Full mixed-strip descent to $R = e/2$]
\label{cor:p12-mixed-strip-descent-to-ehalf}
Let $2a < T_0 < c$, $T < S < T_0$.  Then
\(\ker L_{R,S,T_0}^{\{a,b,2a\}} = \{0\}\)
throughout the mixed strip $e/2 \le R < T$.
\end{corollary}

\begin{proof}
Combine Theorems~\ref{thm:p12-b2a}, \ref{thm:p12-b2c}, and
\ref{thm:p12-b2d}.
\end{proof}

\begin{remark}[Independent inequality inputs]
\label{rem:p12-b2d-inputs}
The proof uses four elementary rational inequalities and no
transcendence input:

\begin{enumerate}
\item[(i)] $p^{2} - q^{2} > 0$, from $(p/q)^{2} = 2^{3/2} > 1$.
\item[(ii)] $2p^{2} - q^{2} > 0$, from $4\sqrt{2} > 1$.
\item[(iii)] $(2 - 3B + B^{2})/(1 + B) < 4/5$ where $B = 2^{-3/2}$,
equivalent to $2809 < 2888$.
\item[(iv)] $\theta > 4/5$ where $\theta = (\log 3 / \log 2)(2/3)^{3/2}$,
via $\log 3 > 2\sqrt{2}$ (equivalent to $9 > 8$) and $50 > 48$.
\end{enumerate}
Only (iv) uses a comparison of $\log 3$ and $\log 2$; that comparison
follows from $3 > 2\sqrt{2}$, i.e. $9 > 8$, elementary.
\end{remark}
